\subsection{Decision-theoretic segmentation summaries}
\label{sec:decision-rules}
Posterior summaries answer different decision questions. The distinction is not a formatting convention: the marginal posterior of $k$ sums over segmentations, while joint MAP scoring maximizes over them.

\begin{proposition}[Three non-equivalent posterior modes]
\label{prop:three-modes}
Define
\begin{align*}
\widehat k_{\rm marg}&\in\arg\max_k\left\{\log p(k)+\log L_{k,n}-\log C_k\right\},\\
\widehat t_{\rm cond}&\in\arg\max_{t\in\mathcal T_{\widehat k_{\rm marg}}}
 \sum_{q=1}^{\widehat k_{\rm marg}}\ell_{t_{q-1},t_q},\\
(\widehat k_{\rm joint},\widehat t_{\rm joint})&\in\arg\max_{k,t\in\mathcal T_k}
\left\{\log p(k)-\log C_k+\sum_{q=1}^{k}\ell_{t_{q-1},t_q}\right\}.
\end{align*}
Here $\ell_{ij}=\log\widetilde A^{(0)}_{ij}$. The first is a marginal-MAP segment count, the second a conditional joint-MAP partition at that count, and the third a global joint-MAP count--partition pair. In general they are distinct decisions.
\end{proposition}
\begin{proof}
\textbf{Step 1 (marginalize partitions).} The evidence for a fixed count is $p(y\mid k)=L_{k,n}/C_k$, so Bayes' rule gives the first maximization after discarding its constant normalizer.
\textbf{Step 2 (optimize within a selected count).} Conditional on $k$, the factors $p(k)$ and $C_k^{-1}$ are constant across admissible partitions. Therefore max-sum backtracking solves the second maximization for $k=\widehat k_{\rm marg}$.
\textbf{Step 3 (optimize jointly).} If $k$ also varies, the complete joint score retains both $\log p(k)$ and $-\log C_k$, giving the third maximization.
\textbf{Step 4 (show the distinction).} For each $k$, the marginal score uses a \emph{sum} $L_{k,n}$ over all paths, whereas the best-path score uses their \emph{maximum}. Counts with many moderately probable paths can have greater marginal support than counts with one sharply preferred path, so no identity between the maximizing counts follows.
\end{proof}

\begin{proposition}[Bayes actions for full-partition and coordinate losses]
\label{prop:loss-bayes-action}
For fixed $k$, 0--1 loss on the entire ordered vector $t=(t_1,\ldots,t_{k-1})$ is minimized by a joint MAP partition. For coordinatewise Hamming loss $\mathcal L(a,t)=\sum_{q=1}^{k-1}\mathbf1\{a_q\ne t_q\}$, constrained to admissible ordered actions $a_1<\cdots<a_{k-1}$, the Bayes action instead solves
\begin{equation}
\max_{1\le a_1<\cdots<a_{k-1}\le n-1}
\sum_{q=1}^{k-1}p(t_q=a_q\mid y,k).
\label{eq:constrained-hamming}
\end{equation}
This constrained optimization has its own additive dynamic program; independent coordinatewise marginal modes need not form an admissible action.
\end{proposition}
\begin{proof}
\textbf{Step 1 (0--1 loss).} Its posterior risk at action $a$ is $1-p(t=a\mid y,k)$, minimized exactly by maximizing the joint posterior mass.
\textbf{Step 2 (Hamming loss).} By linearity of conditional expectation, its posterior risk is $(k-1)-\sum_qp(t_q=a_q\mid y,k)$. Minimizing risk over valid actions is therefore equivalent to Equation~\eqref{eq:constrained-hamming}; dropping the ordering constraint would solve a different optimization problem.
\textbf{Step 3 (dynamic program).} Set $D_{1,h}=p(t_1=h\mid y,k)$ and recursively
\[
D_{q,h}=p(t_q=h\mid y,k)+\max_{g<h}D_{q-1,g}
\]
on valid index ranges, retaining predecessors. Backtracking from the best valid terminal state yields the constrained Bayes action; prefix maxima implement each layer efficiently.
\textbf{Step 4 (explicit counterexample).} For $n=7,k=3$, assign posterior masses $0.17,0.18,0.20,0.20,0.125,0.125$ to the admissible pairs $(1,3),(2,3),(4,5),(4,6),(1,5),(2,6)$, respectively, and zero elsewhere. They sum to one. The marginal mode of $t_1$ is $4$ with probability $0.40$, while that of $t_2$ is $3$ with probability $0.35$, producing the impossible vector $(4,3)$. The joint modes are $(4,5)$ and $(4,6)$, each of mass $0.20$; the constrained Hamming calculation must respect ordering rather than export $(4,3)$.
\end{proof}

The output interface reports each object by its complete name. Posterior-mean segment levels \emph{conditional on} an exported MAP partition are likewise distinct from posterior-mean Bayes curves averaged over the segmentation distribution.
